\documentclass[12pt]{article}
\usepackage[utf8]{inputenc}
\usepackage[T1]{fontenc}
\usepackage{hyperref}
\title{Tshangla Language Resources for K-12 Curriculum Development}
\author{Compiled by Claude Sonnet 5 (claude-sonnet-5) for Andr\'as Kornai}
\date{}
\begin{document}
\maketitle

\section{Introduction}
Tshangla (also called Sharchopikha, ISO 639-3 tsj) is a Sino-Tibetan language of the Bodish branch. It is a lingua franca in eastern Bhutan. The sources I found estimate about 170,000 speakers (157,000 in Bhutan, 11,000 in India, 7,000 in Tibet). They describe it as largely oral, with no established writing system. This differs from the speaker figure in the task description. I found little in the way of educational materials.

\section{Primary Educational Resources}
I found no primers, textbooks or school programmes for Tshangla in these results.

\section{Linguistic and Reference Materials}
\begin{itemize}
\item A Grammar of Tshangla (Brill), described as the first major linguistic description of the language. It covers phonology, lexicon, morphosyntax and discourse. Bookseller listing: \url{https://www.betterworldbooks.com/product/detail/-9789004178274}
\item Tshangla language, Wikipedia (overview): \url{https://wikipedia.com/wiki/Tshangla}
\item A Grammar of Bjokapakha (Selin Grollmann, Brill, 2020), a related language of central Bhutan in the Tshangla branch. It is related material, not Tshangla itself: \url{https://boris.unibe.ch/id/eprint/120832}
\end{itemize}
I did not find a dedicated Tshangla dictionary.

\section{Religious and Community Resources}
None found.

\section{Digital and Audio Resources}
\begin{itemize}
\item Mandala Library (University of Virginia) subject page on Tshangla. I did not verify its contents: \url{https://mandala.library.virginia.edu/subjects/7390/overview/nojs}
\end{itemize}

\section{Curriculum Development Context}
A scholarly grammar exists, but I found no orthography standard, dictionary or classroom materials. A curriculum would need a writing system and basic teaching materials built from the grammar and from community input. In Bhutan the national curriculum does not appear in these results, and I did not check it.

\end{document}